\documentclass{article}
\usepackage{fontspec}
\setmainfont{Microsoft YaHei}
\XeTeXlinebreaklocale "zh"
\XeTeXlinebreakskip = 0pt plus 1pt
\usepackage[margin=2.2cm]{geometry}
\usepackage{booktabs}
\usepackage{enumitem}
\setlist{nosep}
\title{阶段报告（二）：精细手模型、旋转判别力诊断与全量跟踪测试}
\author{}
\date{2026-09-24}
\begin{document}
\maketitle

\section{概要}
本阶段在 Apollo 真实数据集上完成三件事：其一，取得 Barrett Hand 280 的真实网格资产并登记到实测骨架，制成精细手模型；其二，针对"旋转判别力不足"开展定点对照诊断，用实验分离了病因，确认打分方法无辜、病根在模型与照片的两处失配，并以标定常量治愈了其中的系统性部分；其三，在官方"恒定偏差"考卷上运行完整跟踪流程：末端平移误差中位 3.28 厘米，比纯编码器死推算（21.03 厘米）好 84\%——"基线价值"判据通过；"主判据"（1.0 厘米）未通过，与官方融合路线的 0.62 厘米尚有约 5 倍差距，差距已定位到手部朝向的可观测性短板。

\section{一、精细手模型}
模型资产取自公开仓库 jhu-lcsr-attic/bhand\_model：Barrett Hand 280 的十个 STL 网格、URDF 与 SolidWorks 原始零件（论文与数据集宣称的"随附 URDF 与开源代码"实为烂尾，已三路核实，改走公开资产路线）。将真网格登记到数据自带的实测骨架（消息流中的整机 TF 骨骼）：

\begin{itemize}
\item 掌块位置配准：三个指根实测点对网格基座点，残差 2.84 毫米；
\item 框架约定差约 134$^\circ$：两套模型（开源网格的导出系与机器人的关节系）的连杆坐标系约定不同，跨模型直接比对朝向无物理意义——仅记录，配准不使用；
\item 掌块钟差常量 $+25^\circ$：以标定扫描（两轮独立配置：峰位 $+30^\circ$ 与 $+25^\circ$，各自留出帧验证均显著优于零修正）测得的装配角常量，性质同相机-头部标定；
\item 指节网格排除：受上述约定差影响，网格挂实测指节系会自带错误的旋转，部件消融实测"加手指反而掉分"——有据排除，待约定旋转可由关节轴导出后并回。
\end{itemize}

\section{二、旋转判别力定点诊断}
背景：真实照片上的抽查显示平移判别健康（大扰动命中 81\%--92\%）而旋转判别不足（45\%--60\%，其一轴甚至倒挂）。为避免补丁式试错，设计四组对照实验定位病因：

\begin{enumerate}
\item \textbf{零域差对照（关键判别器）}：以模型自身渲染的深度图为"观测"重考——三个旋转轴全部 40/40 满分。结论：打分方法与网格几何具备完全的旋转判别力，病根百分之百在"模型与真实照片的失配"，打分方法无罪。
\item \textbf{分轴检查}：三轴分开判——平移轴健康，一根旋转轴的判别曲线倒挂（扰动越大分数反而越好），指向约 $+25^\circ$ 的系统性装配角错位。
\item \textbf{部件消融}：掌块/手指/前臂分件计分——加手指反而掉分，指姿失配（或约定差所致的错误指姿）实为噪声源。
\item \textbf{标定扫描}（标定帧/留出帧各十帧，判读规则预先写死）：测得掌块钟差 $+25^\circ$--$30^\circ$（两轮互证），倒挂治愈（该轴命中率由 37.5\% 回升至 47.5\%--60\%）；另一轴的"恒定错位"假设被留出帧两次证伪（其残余为非恒定分量，挂账在案）。
\end{enumerate}

诊断终态（20$^\circ$ 扰动逐帧命中率）：平移轴 80\%，两旋转轴 65\% 与 47.5\%--60\%。较初态（含倒挂）全面改善，且每一步改善都由对照实验支撑。

\section{三、全量跟踪测试}
考卷：官方恒定偏差卷（真实视频 + 每关节 $+8.59^\circ$ 篡改读数 + 动捕真值；答案键由两份同录数据的时间戳对齐读数直减测得，与官方说明的 8.6 度吻合）。考核口径：手对相机六自由度位姿误差。

机制：本课题的"偏差状态化退火进化"——未知量为右臂七关节的恒定偏差，候选真姿态由（读数$-$偏差）经辨识运动学模型正解得到，打分为深度残差（精细手点云），每帧一代"选择$\rightarrow$变异"，扰动退火 $2^\circ\!\to\!0.05^\circ$，百帧时序聚合。

过程实录（三次运行，全程记录修复）：首运行暴露关节辨识判据的位置盲缺陷（纯转动件枢轴不动被误判为固定件，读数未被吸收）；次运行暴露轴符号翻转与七对关节抢同一名称的匹配缺陷；第三次运行模型辨识满分（七对各归其名、拟合残差 0.0000 弧度）。

结果（百帧中位，官方口径）见表~\ref{tab:full}。

\begin{table}[h]\centering
\caption{恒定偏差卷全量测试（中位误差）}
\label{tab:full}
\begin{tabular}{lcc}
\toprule
路线 & 平移误差 & 角度误差 \\
\midrule
纯编码器（官方 / 本次复现） & 19.7 / 21.0 cm & 27.0$^\circ$ / 27.2$^\circ$ \\
\textbf{本课题（粒子偏差辨识）} & \textbf{3.28 cm} & 23.5$^\circ$ \\
官方融合（对照目标线） & 0.62 cm & 2.3$^\circ$ \\
纯视觉无先验（官方） & 47.8 cm & 79.8$^\circ$ \\
\bottomrule
\end{tabular}
\end{table}

判据裁定（预先定死）：主判据（平移中位 $\leq$1.0 厘米）3.28 厘米——未通过；基线价值（比死推算低 $\geq$50\%）3.28 对 21.03 厘米（好 84\%）——通过。

补充分析：偏差估计七轴中三轴接近真值（其一 8.38$^\circ$ 对真值 8.59$^\circ$），两轴符号偏离、两轴偏弱；平移大幅改善而角度改善有限（27.2$^\circ$ $\to$ 23.5$^\circ$）——与第二部分的诊断完全同构：手的位置在图像中可观测性好，朝向可观测性弱。死推算基线与官方纯编码器列的吻合（21.0 对 19.7 厘米、27.2$^\circ$ 对 27.0$^\circ$）则验证了整套考核机器的可比性。

\section{四、判据总账}
\begin{table}[h]\centering
\begin{tabular}{lcc}
\toprule
判据组 & 结果 & 备注 \\
\midrule
全量测试-主判据 & 未通过 & 3.28cm（目标 1.0cm） \\
全量测试-基线价值 & \textbf{通过} & 比死推算好 84\% \\
旋转诊断-平移维 & 通过 & 80\% 命中 \\
旋转诊断-旋转维 & 未通过 & 65\% / 47.5--60\%（倒挂已治愈） \\
\bottomrule
\end{tabular}
\end{table}

\section{五、结论与下一步}
结论：视觉修正机制在真实数据上兑现了实质价值（84\% 纠错，判据通过），且全程由对照实验支撑、失败记录在案；精度与官方融合线的差距（约 5 倍）集中于手部朝向的可观测性，该短板的病因、药方与验法均已入档。

下一步三选一：（1）旋转可观测性升级（由关节轴导出网格约定旋合并回指节真网格、加长时序窗口）；（2）换考卷扩展验证（方波偏差卷、遮挡与快速运动卷）；（3）整理中期报告正文。

\end{document}
